\documentclass[11pt]{article}
\usepackage[margin=1in]{geometry}
\usepackage{amsmath,amssymb,amsthm,booktabs,enumitem}
\usepackage[hidelinks]{hyperref}
\newtheorem{definition}{Definition}
\newtheorem{lemma}{Lemma}
\newtheorem{theorem}{Theorem}
\newtheorem{corollary}{Corollary}
\newtheorem{assumption}{Assumption}
\theoremstyle{remark}
\newtheorem{remark}{Remark}

\newcommand{\Eplus}{E^{+}}
\newcommand{\PO}{\xrightarrow{\mathrm{po}}}
\newcommand{\HB}{\xrightarrow{\mathrm{hb}}}
\newcommand{\HBu}{\xrightarrow{\mathrm{hb}^{+}}}
\newcommand{\tid}{\mathrm{tid}}
\newcommand{\loc}{\mathrm{loc}}
\newcommand{\kind}{\mathrm{kind}}
\newcommand{\str}{\mathrm{str}}
\newcommand{\scope}{\mathrm{scope}}
\newcommand{\blk}{\mathrm{blk}}
\newcommand{\warp}{\mathrm{warp}}
\newcommand{\grp}{\mathrm{grp}}
\newcommand{\rel}{\mathrm{rel}}
\newcommand{\acq}{\mathrm{acq}}
\newcommand{\incl}{\mathrm{incl}}
\newcommand{\ms}{\mathrm{ms}}
\newcommand{\conf}{\mathrm{conf}}
\newcommand{\unord}{\mathrm{unord}}
\newcommand{\DR}{\mathrm{DR}}
\newcommand{\SC}{\mathrm{SC}}
\newcommand{\LS}{\mathrm{LS}}
\newcommand{\MS}{\mathrm{MS}}
\newcommand{\seq}{\mathrm{seq}}
\newcommand{\epoch}{\mathrm{epoch}}
\newcommand{\vc}{\mathit{vc}}
\newcommand{\chain}{\rightsquigarrow}
\newcommand{\Chain}{C}
\newcommand{\Pub}{P}
\newcommand{\Fences}{\mathcal{F}}
\newcommand{\Mem}{\mathcal{M}}
\newcommand{\bucket}{\mathrm{key}}

\title{HB race detection: access model, verdicts and proofs\\
\large Revision 4 delta against \texttt{hb\_race\_detection.tex} (v3)}
\author{}
\date{September 2026}

\begin{document}
\maketitle

\noindent\textbf{Scope of this delta.}
v3 classifies an access as \emph{atomic} iff its opcode is a read-modify-write,
and defines a race as a conflicting pair unordered by $\HB$. Under the PTX
memory model (ISA \S8.7) strength is a property of the qualifier, not the
opcode: \texttt{ld}/\texttt{st} with \texttt{.relaxed}/\texttt{.acquire}/%
\texttt{.release}/\texttt{.sc} (and \texttt{.volatile}, which is
\texttt{.relaxed.sys}) are strong, and a conflicting pair that is \emph{morally
strong} is not a data race (\S8.7.1). This delta (i) extends the record
labels with strength and scope, (ii) replaces the single ``race'' verdict by
two verdicts, \emph{data race} and \emph{unordered strong conflict}, exactly
following \S8.7.1, (iii) fence-gates the (ATOM) generator so that $\HB$ is
contained in PTX causality order, (iv) replaces the FastTrack last-write
bookkeeping, which is unsound under (ii), by per-class buckets, and (v)
restates the single-trace theorems and the certificate. Everything not
mentioned (trace model, ghost events, TV1--TV3, the SYNC generator, the
well-formedness monitor, MM1--MM3) is unchanged. Soundness and completeness
keep v3's meaning: \emph{sound} = no missed verdict, \emph{complete} = no
spurious verdict.

\section{Access model}

\begin{definition}[Record labels]\label{def:labels}
Every memory record $e$ carries, in addition to $\tid(e)$, $\mathrm{pc}(e)$,
$\mathrm{addr}(e)$ and $\seq(e)$:
\begin{itemize}[nosep]
\item $\kind(e)\in\{\mathtt{R},\mathtt{W},\mathtt{RMW}\}$;
\item $\str(e)\in\{\mathtt{weak},\mathtt{strong}\}$, with
      $\str(e)=\mathtt{strong}$ whenever $\kind(e)=\mathtt{RMW}$;
\item $\scope(e)\in\{\mathtt{cta},\mathtt{gpu},\mathtt{sys}\}$, defined iff
      $\str(e)=\mathtt{strong}$, ordered $\mathtt{cta}<\mathtt{gpu}<\mathtt{sys}$;
\item for $\kind(e)=\mathtt{RMW}$, flags $\rel(e),\acq(e)\in\{0,1\}$
      (Definition~\ref{def:gate});
\item $\warp(e)$, the warp of $\tid(e)$, and $\grp(e)$, the issue group: two
      records of one warp have the same group iff they were produced by the
      same dynamic instruction of that warp.
\end{itemize}
$\kind$, $\str$, $\scope$, $\rel$, $\acq$ are functions of $\mathrm{pc}(e)$
supplied by the static sidecar (\S\ref{sec:obl}); $\warp,\grp$ come from the
collector. A \emph{write} is a record with $\kind\ne\mathtt{R}$.
\end{definition}

\begin{assumption}[Uniform size, generic proxy, one device]\label{as:size}
All records on a location have the same access size, all are performed via the
generic proxy, and all threads of an execution run on one device. Hence PTX's
``overlap completely'' and ``same proxy'' conditions on moral strength hold
for every pair on a location, and $\mathtt{gpu}$ scope includes every thread.
Mixed-size and cross-proxy accesses are outside the model; the monitor rejects
a trace that has two sizes on one location.
\end{assumption}

\begin{definition}[Scope inclusion, moral strength]\label{def:ms}
For a strong record $e$ and thread $t$,
\[
\incl(e,t)\;:\Longleftrightarrow\;
\scope(e)\in\{\mathtt{gpu},\mathtt{sys}\}\ \vee\
\big(\scope(e)=\mathtt{cta}\wedge \blk(\tid(e))=\blk(t)\big).
\]
Records $a,b$ with $\tid(a)\ne\tid(b)$ are \emph{morally strong},
$\ms(a,b)$, iff $\str(a)=\str(b)=\mathtt{strong}$, $\incl(a,\tid(b))$ and
$\incl(b,\tid(a))$. (PTX \S8.7 also declares same-thread pairs morally
strong; those are $\PO$-ordered and never reach a verdict.)
\end{definition}

Note that $\ms$ is symmetric, is a function of the two labels and the two
block ids, and coincides on RMW pairs with v3's $\mathrm{covers}(r,q)$.

\section{Verdicts}

\begin{definition}[Conflict, unordered, verdicts]\label{def:verdict}
For records $a,b$ with $\seq(a)<\seq(b)$:
\begin{align*}
\conf(a,b)  &:\Longleftrightarrow\ \loc(a)=\loc(b)\ \wedge\ \tid(a)\ne\tid(b)\
             \wedge\ (\kind(a)\ne\mathtt{R}\vee\kind(b)\ne\mathtt{R}),\\
\unord(a,b) &:\Longleftrightarrow\ \neg(a\HB b)\ \wedge\ \neg(b\HB a),\\
\DR(a,b)    &:\Longleftrightarrow\ \conf(a,b)\wedge\unord(a,b)\wedge\neg\ms(a,b)
             \quad\text{(data race, PTX \S8.7.1)},\\
\SC(a,b)    &:\Longleftrightarrow\ \conf(a,b)\wedge\unord(a,b)\wedge\ms(a,b)
             \wedge\neg\big(\kind(a)=\kind(b)=\mathtt{RMW}\big)
             \quad\text{(unordered strong conflict)}.
\end{align*}
A pair is \emph{reportable} if $\DR(a,b)\vee\SC(a,b)$. The \emph{class} of a
reportable pair is $\DR$ or $\SC$.
\end{definition}

\begin{remark}[Why RMW--RMW pairs are excluded from $\SC$]
Two morally strong RMWs on a location are totally ordered by coherence and
each is performed entirely before or after the other (PTX \S8.7.2, atomicity
of RMWs); no update is lost and no torn value is observed. That is the
contract programmers rely on for counters and locks, so reporting every pair of
concurrent \texttt{atomicAdd}s would drown the signal. Every other morally
strong unordered conflict involves a plain strong load or store, for which
coherence gives a total order but no atomicity guarantee: a store may be
overwritten, a load may return either coherence-order value. Those are the
conflicts the tool reports. The exclusion is a reporting choice, not a
semantic claim; flip it by dropping the last conjunct of $\SC$.
\end{remark}

\begin{remark}[What the verdicts do not say]
Neither verdict claims the program is incorrect. $\DR$ says the platform gives
no guarantee about the pair; $\SC$ says the pair's outcome is determined by a
schedule-dependent coherence order. Whether a program tolerates either is a
property of the algorithm, not of the trace, and is outside this document.
\end{remark}

\begin{definition}[Mechanism labels]\label{def:labels2}
Two informational labels refine a reportable pair; they never change its class.
\begin{itemize}[nosep]
\item \emph{Lock-step dependent}:
      $\LS(a,b):\Leftrightarrow \warp(a)=\warp(b)\wedge\grp(a)\ne\grp(b)$.
      Under the pre-Volta execution model a warp issues one instruction at a
      time for all its active lanes, so distinct issue groups of one warp are
      totally ordered by issue order; $\LS$ marks a pair whose only ordering
      was that model, which independent thread scheduling does not provide.
\item \emph{Missing synchronisation}:
      $\MS(a,b):\Leftrightarrow a\HBu b\ \vee\ b\HBu a$, where $\HBu$ is
      Definition~\ref{def:hb} with every $\rel,\acq$ flag set to $1$. $\MS$
      marks a pair that an unqualified atomic chain would order if the chain
      were fenced.
\end{itemize}
\end{definition}

\section{Happens-before with fence-gated atomics}

\begin{definition}[Fence adjacency, static]\label{def:gate}
Let $G$ be the kernel CFG over instruction pcs, $\Mem$ the set of memory-access
pcs, and $\Fences_\sigma$ the set of fence pcs of scope $\ge\sigma$
(\S\ref{sec:obl}, O2). For an RMW pc $r$ with $\scope(r)=\sigma$:
\begin{itemize}[nosep]
\item $\rel(r)=1$ iff every backward path in $G$ from $r$ meets a node of
      $\Fences_\sigma$ before it meets a node of $\Mem\setminus\{r\}$, or
      reaches the entry meeting neither;
\item $\acq(r)=1$ iff every forward path in $G$ from $r$ meets a node of
      $\Fences_\sigma$ before it meets a node of $\Mem\setminus\{r\}$, or
      reaches an exit meeting neither.
\end{itemize}
Equivalently: in $G$ with $\Fences_\sigma$ deleted, no node of
$\Mem\setminus\{r\}$ reaches $r$ (release), respectively is reachable from $r$
(acquire).
\end{definition}

The definition is deliberately stronger than ``a fence dominates $r$''. A
dominating fence hoisted out of a loop orders the accesses of earlier
iterations before the release, but not those of the current iteration; PTX's
release pattern (\S8.9) orders only the accesses that precede the fence in
program order. Adjacency guarantees that every access of the releasing thread
that precedes $r$ in program order also precedes a fence of adequate scope.
The same holds on the acquire side.

\begin{definition}[Release chain]\label{def:chain}
For RMWs on one location, \emph{successor} is as in v3 (next RMW in $\seq$,
which by TV3 is the next in coherence order). $r\chain q$ iff there are RMWs
$r=z_0,z_1,\dots,z_k=q$, $k\ge1$, with $z_{i+1}$ the successor of $z_i$ and
$\ms(z_i,z_{i+1})$ for every $i$.
\end{definition}

\begin{definition}[Happens-before]\label{def:hb}
$\HB$ is the least transitive relation on $\Eplus$ containing
\begin{itemize}[nosep]
\item[(PO)] $a\PO b\Rightarrow a\HB b$;
\item[(SYNC)] as in v3: $\mathrm{arr}_s^B\HB\mathrm{dep}_t^B$ for all
      $s,t\in S(B)$;
\item[(ATOM)] $r\HB q$ for RMWs $r,q$ with $r\chain q$, $\rel(r)=1$ and
      $\acq(q)=1$.
\end{itemize}
$\HBu$ is the same relation with (ATOM) read as $r\chain q$ only.
\end{definition}

\begin{lemma}[Monotonicity]\label{lem:mono}
If $\HB_1\subseteq\HB_2$ on the same trace then every pair reportable under
$\HB_2$ is reportable under $\HB_1$ with the same class.
\end{lemma}
\begin{proof}
$\conf$, $\ms$ and $\kind$ do not depend on $\HB$; $\unord$ is antitone.
\end{proof}

Lemma~\ref{lem:mono} is what makes conservative labelling safe: a sidecar that
fails to recognise a fence, a scope, or a strength only removes (ATOM) edges or
turns $\ms$ false, and can add reports but never hide one.

\begin{lemma}[$\HB\subseteq\seq$]\label{lem:seq}
For every well-formed trace and all $a,b\in\Eplus$: $a\HB b\Rightarrow
\seq(a)<\seq(b)$.
\end{lemma}
\begin{proof}
As in v3 for (PO) and (SYNC). For (ATOM), $r\chain q$ implies $r$ precedes $q$
in coherence order, which is $\seq$ order by TV3.
\end{proof}

\begin{assumption}[Sync locations]\label{as:sync}
If some RMW on location $x$ has $\rel=1$, then every write on $x$ is an RMW.
The monitor checks this; on violation it clears $\rel$ on every RMW of $x$
(Lemma~\ref{lem:mono}: conservative). The reason: a plain strong store to $x$
between $r$ and $q$ in coherence order breaks PTX's observation-order chain
through $x$, and the trace does not record the coherence position of stores
(TV3 covers RMWs only).
\end{assumption}

\section{Algorithm}

Ticks of a thread are its arrivals and its RMWs, in $\PO$; $\epoch(e)$ of a
record $e$ is the number of ticks of $\tid(e)$ that precede $e$ in $\PO$
(a tick's own record counts the ticks strictly before it). $\vc[t]$ is thread
$t$'s vector clock; publishing happens before ticking, as in v3.

\begin{definition}[Per-location state]\label{def:state}
For each location $x$ the algorithm keeps
\begin{itemize}[nosep]
\item buckets $B_x[u,\kappa]$ for every thread $u$ and every
      $\kappa\in\{\mathtt{R},\mathtt{W},\mathtt{RMW}\}\times
      \{\mathtt{weak},\mathtt{strong}\}\times\{\mathtt{cta},\mathtt{gpu},\mathtt{sys},\bot\}$,
      each holding the epoch (and pc) of $u$'s latest record on $x$ with that
      key, where $\bucket(e)=(\kind(e),\str(e),\scope(e))$ and
      $\scope=\bot$ for weak records;
\item a chain clock $\Chain_x$ (a vector clock or $\bot$) and the label
      $(\tid,\scope,\blk)$ of the last RMW on $x$, $\mathrm{last}_x$.
\end{itemize}
\end{definition}

\begin{definition}[Processing a record]\label{def:alg}
On record $a$ of thread $t$ on location $x$, in this order:
\begin{enumerate}[nosep]
\item[(1)] \emph{Chain step} (only if $\kind(a)=\mathtt{RMW}$). If
      $\mathrm{last}_x$ exists and $\neg\ms(\mathrm{last}_x,a)$, set
      $\Chain_x:=\bot$. If $\acq(a)=1$ and $\Chain_x\ne\bot$, set
      $\vc[t]:=\vc[t]\sqcup\Chain_x$.
\item[(2)] \emph{Checks.} For every thread $u\ne t$ and every key $\kappa$
      with an entry $e=B_x[u,\kappa]$: if
      $\big(\kind(e)\ne\mathtt{R}\vee\kind(a)\ne\mathtt{R}\big)$ and
      $\epoch(e)>\vc[t][u]$, report $(e,a)$ with class $\DR$ if
      $\neg\ms(e,a)$, class $\SC$ if $\ms(e,a)$ and not both RMW, and no
      report otherwise.
\item[(3)] \emph{Record.} $B_x[t,\bucket(a)]:=(\epoch(a),\mathrm{pc}(a))$.
\item[(4)] \emph{Publish and tick} (only if $\kind(a)=\mathtt{RMW}$). If
      $\rel(a)=1$, set $\Chain_x:=(\Chain_x\text{ or }\bot)\sqcup\vc[t]$
      ($\bot\sqcup v=v$); set $\mathrm{last}_x:=(t,\scope(a),\blk(t))$; then
      $\vc[t][t]:=\vc[t][t]+1$.
\end{enumerate}
Arrivals, departures and exits are processed as in v3. The $\HBu$ instance
used for the $\MS$ label runs the same procedure with $\rel=\acq=1$
everywhere, on a second copy of the clocks.
\end{definition}

The v3 state (\texttt{last\_write}, \texttt{last\_reads},
\texttt{released[loc]}) is the special case of Definition~\ref{def:state} with
all records weak except RMWs, one write bucket per location, and
$\Chain_x=\texttt{released[loc]}$.

\section{Single-trace results}

\begin{lemma}[Chain clock]\label{lem:chain}
After step (4) of an RMW $z$ on $x$,
\[
\Chain_x=\bigsqcup\{\Pub_r : r\text{ an RMW on }x,\ \rel(r)=1,\ r\chain z\}
\ \sqcup\ (\rel(z)=1\ ?\ \Pub_z : \bot),
\]
where $\Pub_r$ is the value of $\vc[\tid(r)]$ published at $r$ (after $r$'s
own step (1), before its tick).
\end{lemma}
\begin{proof}
Induction over the RMWs on $x$ in $\seq$ order. For the first RMW the chain
set is empty and the claim is step (4). For $z$ with predecessor $p$: by IH
$\Chain_x$ before step (1) of $z$ equals
$\bigsqcup\{\Pub_r:\rel(r),\,r\chain p\}\sqcup(\rel(p)?\Pub_p:\bot)
=\bigsqcup\{\Pub_r:\rel(r),\,r\chain p\text{ or }r=p\}$. By
Definition~\ref{def:chain}, $r\chain z$ iff ($r\chain p$ or $r=p$) and
$\ms(p,z)$; step (1) resets exactly when $\neg\ms(p,z)$, and step (4) adds
$\Pub_z$ iff $\rel(z)$.
\end{proof}

\begin{lemma}[Clock invariant]\label{lem:clock}
For every thread $t$, after processing each event $c$ of $t$, and every thread
$s\ne t$: $\vc[t][s]=\max\{\epoch(o): o\text{ a tick of }s,\ o\HB c\}$
($\max\emptyset=-1$), and $\vc[t][t]$ is the number of ticks of $t$ up to and
including $c$.
\end{lemma}
\begin{proof}
Induction on $\seq$. The cases $c$ a plain record, arrival or departure are
those of v3 (the (PO) and (SYNC) generators are unchanged). Let $c=a$ be an
RMW of $t$ on $x$ and $a^-$ the previous event of $t$. The $\HB$-predecessors
of $a$ are those of $a^-$ together with, for every RMW $r$ on $x$ with
$r\chain a$, $\rel(r)=1$, $\acq(a)=1$: $r$ and every predecessor of $r$. For a
thread $s\ne t$ the maximum tick epoch over the first part is
$\vc[t][s]$ before the step (IH at $a^-$). Over the part contributed by $r$:
if $s=\tid(r)$ it is $\epoch(r)$, since $r$ is itself a tick of $s$ and every
tick of $s$ that $\HB$-precedes $r$ is $\PO$-before it; and
$\Pub_r[s]=\epoch(r)$ because $r$ publishes before it ticks. If $s\ne\tid(r)$
it is $\max\{\epoch(o):o\HB r\}=\vc[\tid(r)][s]$ at $r$ (IH) $=\Pub_r[s]$.
By Lemma~\ref{lem:chain} step (1) joins exactly $\bigsqcup\Pub_r$ over the
qualifying $r$ when $\acq(a)=1$ and the chain is intact, and joins nothing
otherwise, which is also correct because then no (ATOM) edge enters $a$.
Step (4) increments $\vc[t][t]$ once for the tick $a$.
\end{proof}

\begin{corollary}[HB test]\label{cor:test}
For records $e$ of $s$ and $a$ of $t\ne s$ with $\seq(e)<\seq(a)$, evaluated
after step (1) of $a$: $e\HB a\iff\epoch(e)\le\vc[t][s]$.
\end{corollary}
\begin{proof}
$e\HB a$ iff some tick $o$ of $s$ has $e\PO o$ (or $e=o$) and $o\HB a$: the
(SYNC) and (ATOM) edges leaving $s$ originate at arrivals and RMWs, which are
ticks, and conversely $e\PO o\HB a$ gives $e\HB a$. And $e\PO o$ or $e=o$ iff
$\epoch(e)\le\epoch(o)$, because $o$ itself is a tick. Combine with
Lemma~\ref{lem:clock}.
\end{proof}

\begin{theorem}[Complete: no spurious verdict]\label{thm:complete}
Every pair $(e,a)$ reported at step (2) is reportable, and its reported class
is its class.
\end{theorem}
\begin{proof}
$e$ lies in a bucket of $x=\loc(a)$ for a thread $u\ne t$, so $\loc(e)=\loc(a)$
and $\tid(e)\ne\tid(a)$; the kind test gives $\conf$. $\epoch(e)>\vc[t][u]$
gives $\neg(e\HB a)$ by Corollary~\ref{cor:test}, and $\seq(e)<\seq(a)$ gives
$\neg(a\HB e)$ by Lemma~\ref{lem:seq}; so $\unord$. The class is decided by
$\ms(e,a)$ and the two kinds, exactly as in Definition~\ref{def:verdict}.
\end{proof}

\begin{lemma}[Bucket representative]\label{lem:rep}
Let $(e,a)$ be reportable with $\seq(e)<\seq(a)$, and let $e'$ be the entry
of $B_x[\tid(e),\bucket(e)]$ when $a$ is processed. Then $(e',a)$ is
reportable with the same class, and $(e',a)$ is reported at step (2).
\end{lemma}
\begin{proof}
$e'$ is the latest record of $u=\tid(e)$ on $x$ with $\bucket(e')=\bucket(e)$,
so $e=e'$ or $e\PO e'$. If $e'\HB a$ then $e\HB a$, contradicting $\unord(e,a)$;
so $\neg(e'\HB a)$, and $\neg(a\HB e')$ since $\seq(e')<\seq(a)$
(Lemma~\ref{lem:seq}): $\unord(e',a)$. $\kind(e')=\kind(e)$, so
$\conf(e',a)$ iff $\conf(e,a)$, and both RMW-ness and $\ms$ (which depends on
$\str$, $\scope$ and the block of $u$) agree between $e$ and $e'$; the class
is preserved. Step (2) inspects every bucket of every other thread and, by
Corollary~\ref{cor:test}, reports $(e',a)$.
\end{proof}

\begin{theorem}[Sound: no missed verdict]\label{thm:sound}
For every reportable pair $(e,a)$ with $\seq(e)<\seq(a)$ the algorithm
reports, while processing $a$, a pair $(e',a)$ of the same class with
$e'=e$ or $e\PO e'$ and $\mathrm{pc}$ recorded. In particular every location
with a $\DR$ pair gets a $\DR$ report, and every location with an $\SC$ pair
gets an $\SC$ report.
\end{theorem}
\begin{proof}
Lemma~\ref{lem:rep}.
\end{proof}

\begin{remark}[Why v3's last-write bookkeeping is unsound here]
FastTrack keeps one write per location because in a race-free prefix all
writes are $\HB$-totally ordered, so the latest write dominates every earlier
one. Under Definition~\ref{def:verdict} two morally strong writes may be
unordered without any $\DR$; then the latest strong write does not dominate
the earlier one, and a later weak read ordered after the latest write but not
after the earlier one is a missed $\DR$. Concretely: $w_0$ strong by $A$,
$w_1$ strong by $B$, $\unord(w_0,w_1)$; a barrier joins $B$ and $C$; weak read
$r$ by $C$. $(w_0,r)$ is a $\DR$; a last-write check against $w_1$ sees
$w_1\HB r$ and reports nothing. Buckets keyed by $(\kind,\str,\scope)$ per
thread avoid the problem because a representative is always $\PO$-related to
the record it replaces. The FastTrack epoch compression remains valid for the
weak write buckets across threads in a $\DR$-free prefix (any two writes one
of which is weak are $\HB$-comparable), but the proof above does not use it.
\end{remark}

\section{PTX grounding and certificate}

\begin{lemma}[Containment in PTX causality]\label{lem:ptx}
Under Assumptions~\ref{as:size} and \ref{as:sync}, and with the sidecar
obligations of \S\ref{sec:obl} met, every (ATOM) edge of
Definition~\ref{def:hb}, extended by (PO) on both sides, is contained in the
causality order of the PTX memory model (ISA \S8.9; Lustig, Sahasrabuddhe,
Giroux, ASPLOS 2019).
\end{lemma}
\begin{proof}[Proof sketch]
Let $a\PO r$, $r\chain q$, $q\PO b$ with $\rel(r)=\acq(q)=1$ and
$\sigma=\min(\scope(r),\scope(q))$. By fence adjacency there is a fence $F$ of
scope $\ge\sigma$ with $a\PO F\PO r$ and no access between $F$ and $r$, and a
fence $F'$ with $q\PO F'\PO b$ likewise; $(F;r)$ is a release pattern and
$(q;F')$ an acquire pattern on $\loc(r)$. Each link of $r\chain q$ is a
morally strong RMW reading from its coherence predecessor, so $r$ precedes $q$
in observation order. The patterns are morally strong at scope $\sigma$ because
$\ms$ holds along the chain and the fences have scope $\ge\sigma$; hence the
release pattern synchronises with the acquire pattern, and
$a\ \mathrm{po}\ \mathrm{sw}\ \mathrm{po}\ b$ is a base causality edge.
Assumption~\ref{as:sync} excludes the only way the observation chain could be
interrupted, a non-RMW write on $\loc(r)$.
\end{proof}

This discharges obligation (a) of v3's ``Status of MM1--MM3'' remark for the
(ATOM) generator, and retires the certificate-mode workaround that assigned
scope \textsc{None} to unqualified atomics: gating is now part of the
definition.

\begin{theorem}[Per-profile certificate]\label{thm:cert}
Fix $P$, $I$ and a profile $\Pi$ (coherence order of the RMWs, as in v3). If
some $E\in\mathrm{Exec}(P,I,\Pi)$ has an accepted trace with no reportable
pair, then no $E'\in\mathrm{Exec}(P,I,\Pi)$ has a reportable pair.
\end{theorem}
\begin{proof}
v3's first-race prefix argument with ``race'' read as ``reportable pair''.
The argument uses only that in a prefix without such pairs every read has a
unique $\HB$-maximal writer, so that MM2 fixes its reads-from source and
values, control flow and participant sets match event-for-event. With no
reportable pair every two writes on a location that are not both RMWs are
$\HB$-comparable, and RMW--RMW order is fixed by $\Pi$; so the writers of a
read are totally ordered and the maximal one is unique. The remaining steps
(matching, forward and backward $\HB$ preservation, barrier completion via
MM3) do not mention the verdict classes.
\end{proof}

\begin{remark}[No $\DR$-only certificate]
A trace with no $\DR$ but with $\SC$ pairs does not certify: the coherence
order of unordered strong loads and stores is a schedule-dependent input to
the execution, exactly as the RMW order is, and TV3 does not deliver it. A
certificate over ``$\DR=\emptyset$'' would need a profile over all strong
accesses and a collector that records their coherence positions; it is not
claimed.
\end{remark}

\section{Implementation obligations}\label{sec:obl}

\begin{enumerate}[label=O\arabic*.,nosep]
\item \textbf{Sidecar strength and scope.} For every memory pc: $\str$ from
      the presence of \texttt{.STRONG}; $\scope$ from \texttt{.CTA}/%
      \texttt{.GPU}/\texttt{.SYS}; \texttt{volatile} compiles to
      \texttt{.STRONG.SYS} and needs no special case. RMWs: bare
      \texttt{ATOMS} is \texttt{cta}; a bare \texttt{ATOM}/\texttt{ATOMG}
      whose scope cannot be determined is labelled with $\scope=\bot$, which
      makes $\ms$ false (conservative, Lemma~\ref{lem:mono}).
\item \textbf{Fence inventory for the target SM.} Determine empirically
      (compile the PTX qualifiers, disassemble) which SASS sequences realise
      \texttt{fence}/\texttt{membar} at each scope and the compiler's
      \texttt{ld.acquire}/\texttt{st.release}/\texttt{atom.acq\_rel} idioms
      on sm\_86; $\Fences_\sigma$ is that set at scope $\ge\sigma$. The
      inventory is a table in the artifact, not an assumption of the proof.
\item \textbf{Fence adjacency.} Compute $\rel,\acq$ per RMW pc by the
      reachability test of Definition~\ref{def:gate} on $G$ with
      $\Fences_\sigma$ deleted.
\item \textbf{Engine and oracle.} Replace \texttt{last\_write}/%
      \texttt{last\_reads} by the buckets of Definition~\ref{def:state}
      (per thread: $3\times2\times4$ keys, most empty); implement steps
      (1)--(4) including the chain reset; keep the second clock copy for
      $\MS$. Parity test engine\,=\,oracle key-by-key as before.
\item \textbf{Trace fields.} Block id per record (for $\incl$), warp id and
      issue group per record (for $\LS$). Monitor checks: uniform size per
      location (Assumption~\ref{as:size}); sync-location condition
      (Assumption~\ref{as:sync}) with the conservative reaction.
\item \textbf{Litmus additions}, one kernel each, expected verdict fixed
      before running: (i) strong \texttt{sys} store vs strong \texttt{sys}
      load, unordered $\to$ $\SC$ only; (ii) weak vs strong, unordered
      $\to$ $\DR$; (iii) \texttt{cta}-scoped strong pair across blocks $\to$
      $\DR$; (iv) fence--RMW handshake $\to$ ordered; same without the release
      fence $\to$ $\DR$/$\SC$ with $\MS$; same with the fence hoisted before a
      loop that stores in the body $\to$ not ordered; (v) volatile warp tail
      without \texttt{\_\_syncwarp} $\to$ $\SC$ with $\LS$; (vi) the
      last-write counterexample of the remark above $\to$ $\DR(w_0,r)$
      reported; (vii) the existing 33 ScoR litmus and the canary, with any
      changed expectation justified by Definition~\ref{def:verdict}.
\end{enumerate}

\end{document}
